\section{Introducción y repaso del contexto}

Esta clase es la quinta de la asignatura KAIST CS492D ``Design and Learning of Visual Generative Models'', impartida por el profesor Minhyuk Sung. El tema central de la clase es \textbf{Denoising Diffusion Implicit Models (DDIM)}, un trabajo hito para acelerar el proceso de generación de modelos difusos (Song et al., ICLR 2021).

\subsection{Motivación: de DDPM a la muestreo acelerado}

En las clases anteriores, ya hemos estudiado el marco básico de DDPM (Denoising Diffusion Probabilistic Models): mediante un proceso hacia adelante (forward process) que añade ruido gaussiano paso a paso, se convierte la distribución de datos en una distribución gaussiana estándar; y luego entrenando una red neuronal para aprender el proceso inverso (reverse process), se recupera los datos limpios desde el ruido.

\begin{importantbox}{Cuello de botella central de DDPM: velocidad de muestreo}
DDPM generalmente requiere un proceso iterativo de desruido con $T=1000$ pasos para generar una imagen de alta calidad. Cada paso necesita una propagación hacia adelante completa de la red neuronal, lo que hace que la velocidad de generación sea mucho más lenta que modelos de generación en un solo paso como GAN. Cómo reducir el número de pasos de muestreo sin sacrificar la calidad es uno de los desafíos centrales de la investigación sobre modelos difusos.
\end{importantbox}

El conferenciante introduce la motivación de investigación de DDIM contrastando las ventajas y desventajas de tres tipos principales de modelos generativos:

\begin{table}[H]
\centering
\caption{Comparativa de características de los principales modelos generativos}
\begin{tabular}{lccc}
\toprule
\textbf{Modelo} & \textbf{Calidad de muestras} & \textbf{Diversidad} & \textbf{Velocidad de generación} \\
\midrule
GAN & Alta & Baja (colapso de modos) & Rápida (un paso) \\
VAE & Baja & Alta & Rápida (un paso) \\
DDPM & Alta & Alta & Lenta ($\sim$1000 pasos) \\
\bottomrule
\end{tabular}
\end{table}

Los modelos difusos se desempeñan excelentemente tanto en calidad como en diversidad; su único punto débil es la velocidad de generación. DDIM es precisamente el primer avance importante para resolver este problema.

\subsection{Relación entre Score-Based Models y DDPM}

En la parte introductoria, el conferenciante repasa brevemente la relación entre score-based generative models y DDPM. Sus ideas centrales son altamente similares:

\begin{knowledgebox}{Rol central de Score Function (función de puntuación)}
Para la distribución de datos $q(\mathbf{x}_t)$, su score function se define como:
$$
\nabla_{\mathbf{x}_t} \log q(\mathbf{x}_t)
$$
En los modelos basados en puntuaciones (score-based models), entrenamos una red neuronal $s_\theta(\mathbf{x}_t, t)$ para aproximar este score function. Mediante muestreo iterativo con Langevin dynamics, podemos partir de cualquier distribución previa y acercarnos gradualmente a las regiones de alta densidad de la distribución de datos.
\end{knowledgebox}

El flujo básico de los modelos basados en puntuaciones es:
\begin{enumerate}
    \item Añadir diferentes niveles de ruido gaussiano a la distribución de datos original $q(\mathbf{x}_0)$, obteniendo una serie de distribuciones borrosas $q(\mathbf{x}_t)$.
    \item Entrenar una red predictor de puntuaciones (score predictor), minimizando la pérdida L2 para coincidir con el score function en cada nivel de ruido.
    \item Durante el muestreo, comenzar desde un alto nivel de ruido y reducir gradualmente el nivel de ruido utilizando annealed Langevin dynamics.
\end{enumerate}

\begin{warningbox}{Necesidad de ruido multiescala}
Si solo se añade poca cantidad de ruido (varianza pequeña), el score function es casi cero en las regiones de baja densidad lejos de los puntos de datos; la red resulta difícil de aprender con precisión y tampoco puede proporcionar una guía de gradiente efectiva para mover los puntos de muestreo. Si solo se añade mucha cantidad de ruido (varianza grande), aunque cubre un amplio rango, la distribución borrosa difiere mucho de la distribución original. Por tanto, es imperativo utilizar un cronograma de ruido (noise schedule) multiescala que vaya de grande a pequeño: localizar primero de manera gruesa y luego recuperar con detalle —esto es precisamente la idea de diseño central compartida por DDPM y los modelos basados en puntuaciones.
\end{warningbox}

El conferenciante señala que el objetivo de entrenamiento de DDPM (pérdida de predicción de ruido $\|\boldsymbol{\epsilon} - \boldsymbol{\epsilon}_\theta(\mathbf{x}_t, t)\|^2$) es en realidad equivalente al objetivo de score matching. Existe la siguiente relación entre el score function y el predictor de ruido:

$$
\nabla_{\mathbf{x}_t} \log q(\mathbf{x}_t | \mathbf{x}_0) = -\frac{\boldsymbol{\epsilon}}{\sqrt{1 - \bar{\alpha}_t}}
$$

Donde $\boldsymbol{\epsilon}$ es el ruido añadido y $\bar{\alpha}_t$ es el parámetro acumulativo del cronograma de ruido. Por tanto, la red de predicción de ruido $\boldsymbol{\epsilon}_\theta$ está aprendiendo esencialmente una versión escalada del score function.

\subsection{Advertencia: perspectiva de tiempo continuo}

El conferenciante también menciona que cuando extendemos el proceso difuso al dominio de tiempo continuo, el proceso hacia adelante puede describirse mediante una ecuación diferencial estocástica (SDE):
$$
d\mathbf{x} = f(\mathbf{x}, t)\,dt + g(t)\,d\mathbf{w}
$$
Su proceso inverso también tiene una expresión SDE en forma cerrada y sigue dependiendo fundamentalmente del score function. Tanto DDPM como los modelos basados en puntuaciones pueden considerarse diferentes esquemas de discretización de este proceso continuo.

\begin{knowledgebox}{Perspectivas distintas bajo un marco unificado}
DDPM define el proceso hacia adelante e inverso partiendo de un modelo gráfico probabilístico (cadena de Markov); los modelos basados en puntuaciones aprenden el score function partiendo de Langevin dynamics. Ambos son matemáticamente equivalentes, siendo solo diferentes parametrizaciones del mismo marco. Esta perspectiva unificada (Song et al., ICLR 2021, ``Score-Based Generative Modeling through Stochastic Differential Equations'') sienta las bases teóricas para los métodos de aceleración posteriores como DDIM y DPM-Solver.
\end{knowledgebox}

\subsection{Resumen de la sección}

En esta sección se repasó la equivalencia entre DDPM y los modelos basados en puntuaciones, aclarando el cuello de botella de velocidad de muestreo al que enfrentan los modelos difusos. DDPM requiere aproximadamente 1000 pasos iterativos para generar muestras de alta calidad, mientras que DDIM, que se introducirá en esta clase, redefine el proceso hacia adelante (extendiendo desde un proceso de Markov a uno no de Markov), lo que permite reducir el número de pasos a alrededor de 50 sin casi perder calidad.

